\documentclass[11pt,a4paper]{article}

\usepackage[utf8]{inputenc}
\usepackage[T1]{fontenc}
\usepackage[ngerman]{babel}

\usepackage{amsmath,amssymb,amsthm,mathtools}
\usepackage[a4paper,margin=2.5cm]{geometry}
\usepackage{lmodern}
\usepackage{microtype}
\usepackage{hyperref}

\title{Kanonische \(ABC\)-/\(E\)-Zerlegung und energetische Radikalschichtung}
\author{Kollaborative Ausarbeitung}
\date{\today}

\newtheorem{satz}{Satz}
\newtheorem{lemma}[satz]{Lemma}
\newtheorem{korollar}[satz]{Korollar}
\theoremstyle{definition}
\newtheorem{definition}[satz]{Definition}
\newtheorem{beispiel}[satz]{Beispiel}
\theoremstyle{remark}
\newtheorem{bemerkung}[satz]{Bemerkung}

\newcommand{\N}{\mathbb{N}}
\newcommand{\rad}{\operatorname{rad}}
\newcommand{\ggT}{\operatorname{ggT}}
\newcommand{\vp}{v_p}

\begin{document}
	\maketitle
	
	\begin{abstract}
		Für natürliche Zahlen \(n\), die weder durch \(2\) noch durch \(3\) teilbar sind, formulieren wir
		eine kanonische Zerlegung in einen \(ABC\)-Anteil und einen energetischen \(E\)-Anteil.
		Der \(ABC\)-Anteil besteht aus allen Primfaktorpotenzen der Familien \(A,B,C\), der
		\(E\)-Anteil aus allen Primfaktorpotenzen der Familie \(E\). Darüber hinaus besitzt der
		\(E\)-Anteil eine eindeutige geschichtete Radikalzerlegung. Damit ergibt sich eine saubere
		zweistufige Struktur aus Familienzerlegung und energetischer Radikalschichtung.
	\end{abstract}
	
	\section{Familien modulo \(12\)}
	
	\begin{definition}[Die vier Familien]
		Für Primzahlen \(p>3\) definieren wir die vier Familien
		\[
		E:=\{p:\ p\equiv 1 \pmod{12}\},\qquad
		A:=\{p:\ p\equiv 5 \pmod{12}\},
		\]
		\[
		B:=\{p:\ p\equiv 7 \pmod{12}\},\qquad
		C:=\{p:\ p\equiv 11 \pmod{12}\}.
		\]
	\end{definition}
	
	\begin{lemma}[Vollständigkeit der Familien]
		Jede Primzahl \(p>3\) liegt eindeutig in genau einer der vier Familien \(E,A,B,C\).
	\end{lemma}
	
	\begin{proof}
		Eine Primzahl \(p>3\) ist weder durch \(2\) noch durch \(3\) teilbar. Daher gilt
		\[
		p\equiv 1,5,7,\text{ oder }11 \pmod{12}.
		\]
		Dies sind genau die vier obigen Familien. Die Restklassen sind disjunkt, also ist die Zuordnung eindeutig.
	\end{proof}
	
	\section{Kanonische \(ABC\)-/\(E\)-Zerlegung}
	
	\begin{definition}[\(ABC\)-Anteil und energetischer Anteil]
		Sei \(n\in\N\) mit
		\[
		\ggT(n,6)=1.
		\]
		Schreibe die Primfaktorzerlegung von \(n\) als
		\[
		n=\prod_{p\mid n} p^{\vp(n)}.
		\]
		Dann definieren wir den \(ABC\)-Anteil von \(n\) durch
		\[
		n_{ABC}:=
		\prod_{\substack{p^{\vp(n)}\parallel n\\ p\equiv 5,7,11\;(\mathrm{mod}\;12)}} p^{\vp(n)},
		\]
		und den energetischen \(E\)-Anteil durch
		\[
		\mathcal E(n):=
		\prod_{\substack{p^{\vp(n)}\parallel n\\ p\equiv 1\;(\mathrm{mod}\;12)}} p^{\vp(n)}.
		\]
	\end{definition}
	
	\begin{satz}[Kanonische \(ABC\)-/\(E\)-Zerlegung]
		Für jede natürliche Zahl \(n\in\N\) mit
		\[
		\ggT(n,6)=1
		\]
		existiert eine eindeutige Zerlegung
		\[
		n=n_{ABC}\,\mathcal E(n),
		\]
		wobei
		\[
		n_{ABC}
		\]
		genau aus den Primfaktorpotenzen der Familien \(A,B,C\) und
		\[
		\mathcal E(n)
		\]
		genau aus den Primfaktorpotenzen der Familie \(E\) besteht.
	\end{satz}
	
	\begin{proof}
		Nach dem Fundamentalsatz der Arithmetik besitzt \(n\) eine eindeutige Primfaktorzerlegung
		\[
		n=\prod_{p\mid n} p^{\vp(n)}.
		\]
		Nach Lemma 1 liegt jede Primzahl \(p>3\), die \(n\) teilt, eindeutig in genau einer der vier Familien
		\(E,A,B,C\). Daher lässt sich das Produkt eindeutig in zwei Teilprodukte zerlegen:
		eines über alle Primzahlen der Familien \(A,B,C\) und eines über alle Primzahlen der Familie \(E\).
		Dies liefert
		\[
		n=n_{ABC}\,\mathcal E(n).
		\]
		Da sowohl die Primfaktorzerlegung als auch die Familienzuordnung eindeutig sind, ist auch diese
		Zerlegung eindeutig.
	\end{proof}
	
	\begin{korollar}
		Zu jedem \(n\in\N\) mit \(\ggT(n,6)=1\) sind die Zahlen
		\[
		n_{ABC}\qquad\text{und}\qquad \mathcal E(n)
		\]
		kanonisch und eindeutig bestimmt.
	\end{korollar}
	
	\section{Energetische Radikalschichtung}
	
	\begin{definition}[Radikalschichten]
		Sei \(m\in\N\) und
		\[
		m=\prod_{p\mid m} p^{\vp(m)}.
		\]
		Für
		\[
		h:=\max_{p\mid m}\vp(m)
		\]
		definieren wir die \(k\)-te Radikalschicht durch
		\[
		R_k(m):=
		\prod_{\substack{p\mid m\\ \vp(m)\ge k}} p,
		\qquad k=1,\dots,h.
		\]
	\end{definition}
	
	\begin{lemma}[Geschichtete Radikalzerlegung]
		Für jedes \(m\in\N\) gilt
		\[
		m=\prod_{k=1}^{h} R_k(m),
		\qquad
		h=\max_{p\mid m}\vp(m).
		\]
		Außerdem ist jede Schicht \(R_k(m)\) quadratfrei.
	\end{lemma}
	
	\begin{proof}
		Eine Primzahl \(p\mid m\) tritt genau in denjenigen Schichten \(R_k(m)\) auf, für die
		\[
		\vp(m)\ge k
		\]
		gilt. Das sind genau \(\vp(m)\) viele Schichten. Also hat \(p\) im Gesamtprodukt
		\[
		\prod_{k=1}^{h}R_k(m)
		\]
		denselben Exponenten wie in \(m\). Daher stimmen beide Zahlen in ihrer Primfaktorzerlegung überein.
		Da jede Primzahl in einer Schicht höchstens einmal vorkommt, ist jede Schicht quadratfrei.
	\end{proof}
	
	\begin{satz}[Energetische Radikalschichtung]
		Für jede natürliche Zahl \(n\in\N\) mit
		\[
		\ggT(n,6)=1
		\]
		besitzt der energetische Anteil
		\[
		\mathcal E(n)
		\]
		eine eindeutige geschichtete Radikalzerlegung
		\[
		\mathcal E(n)=\prod_{k=1}^{h_E}R_k^{(E)},
		\]
		wobei
		\[
		R_k^{(E)}:=
		\prod_{\substack{p\mid \mathcal E(n)\\ \vp(\mathcal E(n))\ge k}} p
		\]
		quadratfrei ist und nur Primzahlen aus der Familie \(E\) enthält.
	\end{satz}
	
	\begin{proof}
		Dies ist eine unmittelbare Anwendung von Lemma 2 auf \(m=\mathcal E(n)\).
		Da alle Primfaktoren von \(\mathcal E(n)\) per Definition in der Familie \(E\) liegen,
		enthalten auch alle Radikalschichten \(R_k^{(E)}\) ausschließlich \(E\)-Primzahlen.
		Die Eindeutigkeit folgt aus der Eindeutigkeit der Exponenten \(\vp(\mathcal E(n))\).
	\end{proof}
	
	\section{Hauptsatz}
	
	\begin{satz}[Kanonische Zerlegung mit energetischer Radikalschichtung]
		Jede natürliche Zahl \(n\in\N\), die weder durch \(2\) noch durch \(3\) teilbar ist, besitzt eine
		eindeutige Darstellung
		\[
		n=n_{ABC}\,\mathcal E(n),
		\]
		wobei
		\[
		n_{ABC}=
		\prod_{\substack{p^{\vp(n)}\parallel n\\ p\equiv 5,7,11\;(\mathrm{mod}\;12)}} p^{\vp(n)}
		\]
		der \(ABC\)-Anteil und
		\[
		\mathcal E(n)=
		\prod_{\substack{p^{\vp(n)}\parallel n\\ p\equiv 1\;(\mathrm{mod}\;12)}} p^{\vp(n)}
		\]
		der energetische Anteil ist.
		
		Außerdem besitzt der energetische Anteil eine eindeutige geschichtete Radikalzerlegung
		\[
		\mathcal E(n)=\prod_{k=1}^{h_E}R_k^{(E)}.
		\]
		
		Insbesondere ist jede solche Zahl durch ihren \(ABC\)-Anteil und die Radikalschichten ihres
		energetischen Anteils eindeutig bestimmt.
	\end{satz}
	
	\begin{proof}
		Die Existenz und Eindeutigkeit der Zerlegung
		\[
		n=n_{ABC}\,\mathcal E(n)
		\]
		ist Satz 1. Die eindeutige geschichtete Radikalzerlegung von \(\mathcal E(n)\) ist Satz 2.
		Daher ist \(n\) eindeutig durch diese beiden Komponenten bestimmt.
	\end{proof}
	
	\section{Beispiele}
	
	\begin{beispiel}
		Für
		\[
		n=5\cdot 7\cdot 11\cdot 13^2\cdot 37
		\]
		gilt
		\[
		n_{ABC}=5\cdot 7\cdot 11,
		\qquad
		\mathcal E(n)=13^2\cdot 37.
		\]
		Die energetische Radikalschichtung ist
		\[
		R_1^{(E)}=13\cdot 37,
		\qquad
		R_2^{(E)}=13.
		\]
		Also
		\[
		\mathcal E(n)=(13\cdot 37)\cdot 13.
		\]
	\end{beispiel}
	
	\begin{beispiel}
		Für
		\[
		n=13\cdot 61^3
		\]
		ist
		\[
		n_{ABC}=1,
		\qquad
		\mathcal E(n)=13\cdot 61^3.
		\]
		Die energetischen Schichten sind
		\[
		R_1^{(E)}=13\cdot 61,\qquad
		R_2^{(E)}=61,\qquad
		R_3^{(E)}=61.
		\]
	\end{beispiel}
	
	\begin{beispiel}
		Für
		\[
		n=5^2\cdot 7\cdot 29\cdot 31
		\]
		liegen alle Primfaktoren in den Familien \(A\) oder \(B\). Daher ist
		\[
		n_{ABC}=n,
		\qquad
		\mathcal E(n)=1.
		\]
		In diesem Fall ist der energetische Anteil trivial.
	\end{beispiel}
	
	\section{Schlussbemerkung}
	
	Die hier formulierte Struktur trennt zwei Ebenen:
	
	\begin{enumerate}
		\item die globale Zerlegung einer zu \(6\) teilerfremden Zahl in einen \(ABC\)-Anteil und einen
		energetischen \(E\)-Anteil,
		\item die innere geschichtete Radikalzerlegung des energetischen Anteils.
	\end{enumerate}
	
	Damit erhält man keine Darstellung als Produkt \emph{einer einzigen} \(ABC\)-Primzahl und eines
	Restes, sondern eine kanonische und vollständige Familienzerlegung, bei der der \(E\)-Anteil
	zusätzlich eine eindeutige energetische Radikalschichtung trägt.
	
\end{document}